% !TEX program = xelatex
% ============================================================
%  第 4 章 Convex Optimization — 基于模板的结构化笔记
% ============================================================

\documentclass[11pt,a4paper]{article}

% ------------------------------------------------------------
%  字体与页面
% ------------------------------------------------------------
\IfFileExists{geometry.sty}{%
  \usepackage[margin=2.5cm]{geometry}
}{%
  \PackageWarning{geometry}{未找到 geometry 宏包，已使用默认 A4 页边距。}%
}
\usepackage{ctex}
\usepackage{microtype}
\linespread{1.08}

% ------------------------------------------------------------
%  工具宏包
% ------------------------------------------------------------
\usepackage{amsmath,amssymb,amsthm}
\usepackage{mathtools}
\usepackage{graphicx}
\usepackage{enumitem}
\usepackage[dvipsnames]{xcolor}
\usepackage[most]{tcolorbox}
\usepackage{titlesec}
\setlength{\parskip}{0.35em}
\setlength{\parindent}{2em}
\setlist[itemize]{itemsep=0.2em, topsep=0.25em}
\setlist[enumerate]{itemsep=0.2em, topsep=0.25em}

% ------------------------------------------------------------
%  超链接
% ------------------------------------------------------------
\definecolor{mydarkblue}{rgb}{0.0,0.08,0.45}
\usepackage{hyperref}
\hypersetup{
  colorlinks = true,
  linkcolor  = mydarkblue,
  citecolor  = mydarkblue,
  urlcolor   = mydarkblue
}

% ------------------------------------------------------------
%  章节样式
% ------------------------------------------------------------
\titleformat{\section}
  {\large\bfseries\color{mydarkblue}}{\thesection}{0.5em}{}[\titlerule]
\titlespacing*{\section}{0pt}{1.5ex}{1ex}
\titleformat{\subsection}
  {\normalsize\bfseries}{\thesubsection}{0.5em}{}
\titlespacing*{\subsection}{0pt}{1ex}{0.5ex}

% ------------------------------------------------------------
%  顶栏 + 横向导航
% ------------------------------------------------------------
\makeatletter
\newcommand{\makeChapterFourTitle}{%
  \begin{tcolorbox}[
    colback=mydarkblue!5, colframe=mydarkblue!10,
    arc=3pt, boxrule=0.5pt,
    left=5pt, right=5pt, top=5pt, bottom=5pt
  ]
    {\large\bfseries\@title}\hfill{\small\@author}\\[-2pt]
    {\tiny\color{gray}\@date\hfill Chapter 4 · Convex optimization notes}
  \end{tcolorbox}
  \vspace{-0.6em}
}
\makeatother

\newtcolorbox{tabNavBoxSmall}[1][]{
  enhanced,
  colframe=white, colback=gray!5, coltext=black!80,
  fontupper=\sffamily\small\bfseries,
  arc=1.5mm, boxrule=0pt, bottomrule=2pt,
  left=2mm, right=2mm, top=1.5mm, bottom=1.5mm,
  #1
}

\newcommand{\horizontalTOCChFour}{%
  \begin{center}
  \begin{tcbraster}[
    raster columns=3, raster width=\linewidth,
    raster column skip=1mm, raster equal height
  ]
    \begin{tabNavBoxSmall}
      \hyperref[sec:ch4-basic-form]{\ 基本形式与定义}
    \end{tabNavBoxSmall}
    \begin{tabNavBoxSmall}
      \hyperref[sec:ch4-feasibility]{\ 可行性与最优性}
    \end{tabNavBoxSmall}
    \begin{tabNavBoxSmall}
      \hyperref[sec:ch4-transform]{\ 等价转化}
    \end{tabNavBoxSmall}
  \end{tcbraster}
  \vspace{-2.5pt}%
  \noindent\textcolor{gray!30}{\rule{\linewidth}{1pt}}
  \end{center}
  \vspace{1.2em}
}

% ------------------------------------------------------------
%  定理环境
% ------------------------------------------------------------
\tcbset{
  mytheo/.style={
    enhanced, breakable,
    colback=#1!5!white, colframe=#1!70!black,
    coltitle=white, fonttitle=\bfseries\small,
    boxrule=0.6pt, arc=3pt,
    left=5pt, right=5pt, top=10pt, bottom=5pt,
    attach boxed title to top left={yshift*=-\tcboxedtitleheight/2, xshift=2mm},
    boxed title style={
      colback=#1!80!black, colframe=#1!80!black,
      boxrule=0pt, arc=2pt
    }
  }
}
\newtcbtheorem[auto counter]            {theorem}    {定理}{mytheo=WildStrawberry}{thm}
\newtcbtheorem[use counter from=theorem]{definition} {定义}{mytheo=NavyBlue}{def}
\newtcbtheorem[use counter from=theorem]{example}    {例子}{mytheo=YellowOrange}{ex}
\newtcbtheorem[use counter from=theorem]{remark}     {注解}{mytheo=Periwinkle}{rem}

\title{第 4 章\quad Convex Optimization 笔记}
\author{手写笔记转 LaTeX（模板化整理）}
\date{\today}

\begin{document}

\makeChapterFourTitle
\horizontalTOCChFour
\begin{center}
  \includegraphics[width=\linewidth]{\detokenize{ChatGPT Image 2026年4月27日 19_35_30.png}}
\end{center}
\vspace{0.5em}

\section{基本形式与定义}\label{sec:ch4-basic-form}

\begin{definition}{一般优化问题}{general-opt}
一般优化问题可写为
\[
\begin{aligned}
\text{minimize}\quad & f_0(x) \\
\text{subject to}\quad & f_i(x)\le 0,\quad i=1,\dots,m,\\
& h_j(x)=0,\quad j=1,\dots,p.
\end{aligned}
\]
其中 $x\in\mathbb{R}^n$，$f_0$ 为目标函数，$f_i$ 为不等式约束函数，$h_j$ 为等式约束函数。
\end{definition}

\begin{definition}{问题域与可行集}{domain-feasible}
问题的定义域记为
\[
D=\bigcap_{i=0}^m \operatorname{dom} f_i \;\cap\; \bigcap_{j=1}^p \operatorname{dom} h_j.
\]
若 $x\in D$ 且满足
\[
f_i(x)\le 0\;(i=1,\dots,m),\qquad h_j(x)=0\;(j=1,\dots,p),
\]
则称 $x$ 为可行点。所有可行点构成可行集
\[
\mathcal X_f=\{x\in D\mid f_i(x)\le0,\ h_j(x)=0\}.
\]
\end{definition}

\section{可行性与最优性}\label{sec:ch4-feasibility}

\begin{definition}{最优值与最优解集}{optimal-value-solution}
最优值定义为
\[
p^\star=\inf\{f_0(x)\mid x\in\mathcal X_f\}.
\]
若 $\mathcal X_f=\varnothing$，则约定 $p^\star=+\infty$。若可行点 $x^\star$ 满足 $f_0(x^\star)=p^\star$，则称 $x^\star$ 为最优解。最优解集记为
\[
\mathcal X_{\mathrm{opt}}
=\{x\in\mathcal X_f\mid f_0(x)=p^\star\}.
\]
\end{definition}

\begin{definition}{$\varepsilon$-次优解集}{$\varepsilon$-suboptimal}
给定 $\varepsilon\ge0$，$\varepsilon$-次优解集定义为
\[
\mathcal X_\varepsilon
=\{x\in\mathcal X_f\mid f_0(x)\le p^\star+\varepsilon\}.
\]
当 $\varepsilon=0$ 时，$\mathcal X_\varepsilon=\mathcal X_{\mathrm{opt}}$。
\end{definition}

\begin{definition}{局部最优解}{local-opt}
可行点 $x$ 称为局部最优解，若存在 $R>0$ 使得
\[
f_0(x)=\inf\Bigl\{f_0(z)\,\Bigm|\,z\in\mathcal X_f,\ \|z-x\|_2\le R\Bigr\}.
\]
\end{definition}

\begin{definition}{可行性问题}{feasibility-problem}
可行性问题只关心约束是否可满足：
\[
\begin{aligned}
\text{find}\quad & x\\
\text{subject to}\quad & f_i(x)\le0,\quad i=1,\dots,m,\\
& h_j(x)=0,\quad j=1,\dots,p.
\end{aligned}
\]
\end{definition}

\section{等价转化方式}\label{sec:ch4-transform}

\subsection{Box constraints 的标准化}
\begin{example}{区间约束到不等式约束}{box-constraints}
问题
\[
\begin{aligned}
\text{minimize}\quad & f_0(x)\\
\text{subject to}\quad & l_i\le x_i\le u_i,\quad i=1,\dots,n
\end{aligned}
\]
可等价写为
\[
\begin{aligned}
\text{minimize}\quad & f_0(x)\\
\text{subject to}\quad & l_i-x_i\le0,\quad i=1,\dots,n,\\
& x_i-u_i\le0,\quad i=1,\dots,n.
\end{aligned}
\]
\end{example}

\subsection{变量代换（change of variables）}
\begin{definition}{通过一一映射重参数化}{change-of-variables}
设 $\phi:\mathbb R^n\to\mathbb R^n$ 在其定义域上为一一映射，且 $\phi(\operatorname{dom}\phi)\supseteq D$。令
\[
\tilde f_i(z)=f_i(\phi(z)),\quad i=0,1,\dots,m,\qquad
\tilde h_j(z)=h_j(\phi(z)),\quad j=1,\dots,p.
\]
则原问题可等价写为
\[
\begin{aligned}
\text{minimize}\quad & \tilde f_0(z)\\
\text{subject to}\quad & \tilde f_i(z)\le0,\quad i=1,\dots,m,\\
& \tilde h_j(z)=0,\quad j=1,\dots,p.
\end{aligned}
\]
\end{definition}

\subsection{单调函数复合变换}
\begin{definition}{目标与约束的等价复合}{monotone-transform}
设 $\psi_0:\mathbb R\to\mathbb R$ 为单调递增函数；对 $i=1,\dots,m$，$\psi_i:\mathbb R\to\mathbb R$ 满足
\[
\psi_i(u)\le0\iff u\le0;
\]
对 $j=1,\dots,p$，$\varphi_j:\mathbb R\to\mathbb R$ 满足
\[
\varphi_j(u)=0\iff u=0.
\]
定义
\[
\tilde f_0(x)=\psi_0(f_0(x)),\quad
\tilde f_i(x)=\psi_i(f_i(x)),\ i=1,\dots,m,\quad
\tilde h_j(x)=\varphi_j(h_j(x)),\ j=1,\dots,p.
\]
则等价问题为
\[
\begin{aligned}
\text{minimize}\quad & \tilde f_0(x)\\
\text{subject to}\quad & \tilde f_i(x)\le0,\quad i=1,\dots,m,\\
& \tilde h_j(x)=0,\quad j=1,\dots,p.
\end{aligned}
\]
\end{definition}

\subsection{松弛变量（slack variables）}
\begin{definition}{不等式到等式的松弛化}{slack-variable}
原问题
\[
\begin{aligned}
\text{minimize}\quad & f_0(x)\\
\text{subject to}\quad & f_i(x)\le0,\quad i=1,\dots,m,\\
& h_j(x)=0,\quad j=1,\dots,p
\end{aligned}
\]
可等价写为
\[
\begin{aligned}
\text{minimize}\quad & f_0(x)\\
\text{subject to}\quad & s_i\ge0,\quad i=1,\dots,m,\\
& f_i(x)+s_i=0,\quad i=1,\dots,m,\\
& h_j(x)=0,\quad j=1,\dots,p.
\end{aligned}
\]
其中 $s=(s_1,\dots,s_m)$ 为松弛变量。
\end{definition}

\subsection{对部分变量先优化（partial minimization）}
\begin{definition}{先对部分变量取下确界}{optimizing-over-some-variables}
对二元函数 $f(x,y)$，定义
\[
\tilde f(x)=\inf_y f(x,y),
\]
则总有
\[
\inf_{x,y} f(x,y)=\inf_x \tilde f(x).
\]
\end{definition}

\begin{example}{分块变量下的等价化简}{partial-min-example}
设 $x=\begin{pmatrix}x_1\\x_2\end{pmatrix}$，$x_1\in\mathbb R^{n_1}$，$x_2\in\mathbb R^{n_2}$。考虑
\[
\begin{aligned}
\text{minimize}\quad & f_0(x_1,x_2)\\
\text{subject to}\quad & f_i(x_1,x_2)\le0,\quad i=1,\dots,m_1,\\
& \hat f_j(x_2)\le0,\quad j=1,\dots,m_2.
\end{aligned}
\]
定义可行集合
\[
\hat{\mathcal Z}=\{z\in\mathbb R^{n_2}\mid \hat f_j(z)\le0,\ j=1,\dots,m_2\},
\]
并令
\[
\tilde f_0(x_1)=\inf_{z\in\hat{\mathcal Z}} f_0(x_1,z),\qquad
\tilde f_i(x_1)=\inf_{z\in\hat{\mathcal Z}} f_i(x_1,z),\ i=1,\dots,m_1.
\]
则可得到仅关于 $x_1$ 的等价（或更紧）表述：
\[
\begin{aligned}
\text{minimize}\quad & \tilde f_0(x_1)\\
\text{subject to}\quad & \tilde f_i(x_1)\le0,\quad i=1,\dots,m_1.
\end{aligned}
\]
\end{example}

\subsection{Epigraph 形式}
\begin{definition}{目标最小化到上图形式}{epigraph-form}
原问题
\[
\begin{aligned}
\text{minimize}\quad & f_0(x)\\
\text{subject to}\quad & f_i(x)\le0,\quad i=1,\dots,m,\\
& h_j(x)=0,\quad j=1,\dots,p
\end{aligned}
\]
可等价改写为
\[
\begin{aligned}
\text{minimize}\quad & t\\
\text{subject to}\quad & f_0(x)-t\le0,\\
& f_i(x)\le0,\quad i=1,\dots,m,\\
& h_j(x)=0,\quad j=1,\dots,p.
\end{aligned}
\]
\end{definition}

\begin{remark}{为什么“等价”}{equivalence-note}
上述转化保持可行域不变，且在目标函数部分保持解的排序关系（特别是 $\psi_0$ 单调递增时）。因此原问题与转化后问题具有相同的最优解集合。
\end{remark}

\section{Convex Optimization Problem}\label{sec:ch4-convex-opt}

\begin{definition}{凸优化问题的标准形式}{convex-opt-standard}
考虑问题
\[
\begin{aligned}
\text{minimize}\quad & f_0(x)\\
\text{subject to}\quad & f_i(x)\le0,\quad i=1,\dots,m,\\
& a_i^Tx=b_i,\quad i=1,\dots,p.
\end{aligned}
\]
若 $f_i\ (i=0,1,\dots,m)$ 均为凸函数，且等式约束为仿射约束，则该问题是凸优化问题。
\end{definition}

\begin{theorem}{凸优化中的局部最优与全局最优}{local-global-convex}
在凸优化问题中，任一局部最优解都是全局最优解。
\end{theorem}

\begin{definition}{可微目标的一阶最优性条件}{first-order-optimality}
设 $\mathcal X_f$ 为可行集。若 $f_0$ 可微且问题为凸优化，则
\[
x^\star \text{ 为最优解}
\iff
x^\star\in\mathcal X_f,\ 
\nabla f_0(x^\star)^T(y-x^\star)\ge0,\ \forall y\in\mathcal X_f.
\]
\end{definition}

\begin{example}{仅含等式约束时的一阶条件}{equality-constrained-stationarity}
考虑
\[
\min f_0(x)\quad \text{s.t.}\quad Ax=b,\qquad \operatorname{dom}f_0=\mathbb R^n.
\]
由一阶条件可得
\[
\nabla f_0(x^\star)^T(y-x^\star)\ge0,\quad \forall y\ \text{s.t.}\ Ay=b.
\]
令 $y=x^\star+v$，则 $v\in\mathcal N(A)=\{d\mid Ad=0\}$，于是
\[
\nabla f_0(x^\star)^Tv\ge0,\quad \forall v\in\mathcal N(A).
\]
又因 $\mathcal N(A)$ 为线性子空间，$-v\in\mathcal N(A)$，故亦有
\[
\nabla f_0(x^\star)^Tv\le0.
\]
综合得
\[
\nabla f_0(x^\star)^Tv=0,\quad \forall v\in\mathcal N(A),
\]
即 $\nabla f_0(x^\star)\perp\mathcal N(A)$，等价于
\[
\nabla f_0(x^\star)\in\mathcal R(A^T).
\]
因此存在 $\lambda\in\mathbb R^p$ 使得
\[
\nabla f_0(x^\star)+A^T\lambda=0.
\]
\end{example}

\section{Quasi-convex Optimization}\label{sec:ch4-quasi-convex}

\begin{definition}{拟凸优化问题}{quasi-convex-problem}
考虑
\[
\begin{aligned}
\text{minimize}\quad & f_0(x)\\
\text{subject to}\quad & f_i(x)\le0,\quad i=1,\dots,m,\\
& Ax=b.
\end{aligned}
\]
其中 $f_0$ 为 quasi-convex 函数，$f_i$ 为凸函数，$Ax=b$ 为仿射约束。
\end{definition}

\begin{remark}{可微情形下的一个充分条件}{quasi-convex-sufficient}
设 $f_0$ 可微，$\mathcal X_f=\{x\mid f_i(x)\le0,\ h_j(x)=0\}$。笔记中常用的一个充分条件是
\[
\nabla f_0(x^\star)^T(y-x^\star)>0,\quad
\forall y\in\mathcal X_f\setminus\{x^\star\}.
\]
该条件是充分条件，不是充要条件。
\end{remark}

\section{典型 Convex Problems}\label{sec:ch4-typical-problems}

\subsection{线性规划（Linear Programming）}
\begin{definition}{线性规划标准型}{linear-programming}
线性规划可写为
\[
\begin{array}{ll}
\text{minimize} & c^Tx+d\\
\text{subject to} & Gx\preceq h\\
& Ax=b,
\end{array}
\]
其中 $c\in\mathbb R^n,\ d\in\mathbb R,\ G\in\mathbb R^{m\times n},\ h\in\mathbb R^m,\ A\in\mathbb R^{p\times n},\ b\in\mathbb R^p$。目标函数和约束函数均为线性（仿射）函数。
\end{definition}

\begin{remark}{广义不等式记号约定}{generalized-inequality-convention}
本节中，向量约束的 $\preceq,\succeq$ 表示相对于非负正交锥 $\mathbb R_+^m$ 的广义不等式（即逐元素不等式）；矩阵约束的 $\preceq,\succeq$ 表示相对于半正定锥 $\mathbf S_+^n$ 的 Loewner 序。
\end{remark}

\begin{example}{LP 的两类常见等价变换}{lp-equivalent-transform}
\textbf{(a) 不等式转等式：}
\[
\begin{array}{ll}
\text{minimize} & c^Tx+d\\
\text{subject to} & Gx+s=h\\
& Ax=b\\
& s\succeq0.
\end{array}
\]

\textbf{(b) 自由变量分解：}
令 $x=x^+-x^-$，其中 $x^+\succeq0,\ x^-\succeq0$，则
\[
\begin{array}{ll}
\text{minimize} & c^Tx^+-c^Tx^-+d\\
\text{subject to} & Gx^+-Gx^-+s=h\\
& Ax^+-Ax^-=b\\
& s\succeq0,\ x^+\succeq0,\ x^-\succeq0.
\end{array}
\]
\end{example}

\begin{definition}{LP 的另一种标准形式}{lp-alternative-standard-form}
线性规划也常写为
\[
\begin{aligned}
\text{minimize}\quad & c^Tx\\
\text{subject to}\quad & Ax=b,\\
& x\succeq0.
\end{aligned}
\]
\end{definition}

\begin{example}{线性分数规划（Linear Fractional Programming, LFP）}{lfp-to-lp}
考虑线性分数规划
\[
\begin{aligned}
(\mathrm P0)\quad
\text{minimize}\quad & f_0(x)\\
\text{subject to}\quad & Gx\preceq h,\\
& Ax=b,
\end{aligned}
\]
其中
\[
f_0(x)=\frac{c^Tx+d}{e^Tx+f},\qquad
\operatorname{dom}f_0=\{x\mid e^Tx+f>0\}.
\]
令
\[
y=\frac{x}{e^Tx+f},\qquad z=\frac{1}{e^Tx+f},
\]
则可得到等价线性规划
\[
\begin{aligned}
(\mathrm P1)\quad
\text{minimize}\quad & c^Ty+dz\\
\text{subject to}\quad & Gy-hz\preceq0,\\
& Ay-bz=0,\\
& e^Ty+fz=1,\\
& z\ge0.
\end{aligned}
\]
\end{example}

\subsection{二次优化（Quadratic Optimization）}
\begin{definition}{二次规划（QP）}{qp-standard-form}
二次规划问题可写为
\[
\begin{aligned}
\text{minimize}\quad & \frac12x^TPx+q^Tx+r\\
\text{subject to}\quad & Gx\preceq h,\\
& Ax=b,
\end{aligned}
\]
其中 $P\in\mathbf S_+^n,\ G\in\mathbb R^{m\times n},\ A\in\mathbb R^{p\times n}$。
\end{definition}

\begin{definition}{QCQP（Quadratically Constrained Quadratic Program）}{qcqp-standard-form}
\[
\begin{aligned}
\text{minimize}\quad & \frac12x^TP_0x+q_0^Tx+r_0\\
\text{subject to}\quad & \frac12x^TP_ix+q_i^Tx+r_i\le0,\quad i=1,\dots,m,\\
& Ax=b.
\end{aligned}
\]
\end{definition}

\subsection{半正定规划（Semidefinite Programming, SDP）}
\begin{definition}{SDP 标准形式}{sdp-standard-form}
半正定规划的标准形式为
\[
\begin{aligned}
\text{minimize}\quad & \operatorname{tr}(CX)\\
\text{subject to}\quad & \operatorname{tr}(A_iX)=b_i,\quad i=1,\dots,p,\\
& X\succeq0,
\end{aligned}
\]
其中 $X\in\mathbf S_+^n,\ C\in\mathbf S^n,\ A_i\in\mathbf S^n,\ b_i\in\mathbb R$。
\end{definition}

\begin{remark}{迹内积表示}{trace-inner-product}
对于对称矩阵变量，目标可看作矩阵内积：
\[
\operatorname{tr}(CX)=\sum_{i,j}C_{ij}X_{ij}.
\]
\end{remark}

\begin{example}{SDP 的对角特例}{sdp-diagonal-special-case}
设 $X\in\mathbf S^n$ 为对角矩阵，并记
\[
x=\operatorname{diag}(X)\in\mathbb R^n,
\]
即 $\operatorname{diag}(X)$ 表示由对角矩阵 $X$ 的对角元素构成的列向量。考虑
\[
\begin{aligned}
\text{minimize}\quad & \operatorname{tr}(CX)\\
\text{subject to}\quad & \operatorname{tr}(A_iX)=b_i,\quad i=1,\dots,p,\\
& X\succeq0,\quad X\ \text{is diagonal}.
\end{aligned}
\]
令 $c=\operatorname{diag}(C)$、$a_i=\operatorname{diag}(A_i)$，则上式等价为
\[
\begin{aligned}
\text{minimize}\quad & c^Tx\\
\text{subject to}\quad & a_i^Tx=b_i,\quad i=1,\dots,p,\\
& x\ge0.
\end{aligned}
\]
因此该对角情形退化为线性规划（LP）。
\end{example}

\begin{definition}{SDP 的常见等价写法}{sdp-equivalent-form}
另一种常见写法为
\[
\begin{aligned}
\text{minimize}\quad & c^Tx\\
\text{subject to}\quad & \sum_{i=1}^n x_iF_i+G\preceq0,\\
& Ax=b,
\end{aligned}
\]
其中 $F_i,G\in\mathbf S^k,\ A\in\mathbb R^{p\times n}$。
\end{definition}

\begin{example}{谱范数（最大奇异值）最小化的 SDP 形式}{spectral-norm-sdp}
设
\[
A(x)=A_0+x_1A_1+\cdots+x_nA_n,\qquad
A_i\in\mathbb R^{p\times q},\ x\in\mathbb R^n.
\]
考虑问题
\[
\min_{x}\ \|A(x)\|_2.
\]
引入标量变量 $t$ 后，可写为等价形式
\[
\begin{aligned}
\text{minimize}\quad & t\\
\text{subject to}\quad &
\begin{bmatrix}
tI_p & A(x)\\
A(x)^T & tI_q
\end{bmatrix}\succeq0.
\end{aligned}
\]
这给出了一个标准的 LMI 约束 SDP。
\end{example}

\begin{remark}{其他常见等价写法}{spectral-norm-sdp-forms}
同一约束还可写成（在 $t\ge0$ 下）
\[
A(x)^TA(x)-t^2I_q\preceq0,
\]
以及对称块矩阵的等价形式
\[
\begin{bmatrix}
tI_q & A(x)^T\\
A(x) & tI_p
\end{bmatrix}\succeq0.
\]
这些形式可由 Schur 补与块矩阵恒等关系互相转换。
\end{remark}

\subsection{二阶锥规划（Second-Order Cone Program, SOCP）}
\begin{definition}{SOCP 标准形式}{socp-standard-form}
二阶锥规划可写为
\[
\begin{aligned}
\text{minimize}\quad & f^Tx\\
\text{subject to}\quad & \|A_ix+b_i\|_2\le c_i^Tx+d_i,\quad i=1,\dots,m,\\
& Fx=g.
\end{aligned}
\]
其中每个
\[
\|A_ix+b_i\|_2\le c_i^Tx+d_i
\]
称为二阶锥约束（second-order cone constraint）。
\[
\text{等价地，}\ (c_i^Tx+d_i,\ A_ix+b_i)\in\mathcal Q^{k_i+1},
\]
其中 $k_i=\dim(A_ix+b_i)$。
即相对于二阶锥 $\mathcal Q^{k_i+1}$ 的广义不等式约束。
\end{definition}

\end{document}
